\Needspace{8\baselineskip}
\Article{194}
Since the forms we called associated above (Article 187, 6) are always improperly equivalent (Article 159), it is clear that if the reduced forms $F$, $f$ are improperly equivalent, and $F$ has the associated form $G$, then $f$, $G$ are properly equivalent, so $G$ is contained in the period of $f$. Thus if the forms $F$, $f$ are both properly and improperly equivalent, both $F$ and $G$ must occur in the period of $f$. This period is therefore associated with itself and contains two ambiguous forms (Article 187, 7). This gives a fine confirmation of the theorem in Article 165, which already assured us that there is an ambiguous form equivalent to $F$, $f$.

\Needspace{8\baselineskip}
\Article{195}
\textsc{Problem.} \textit{Given any two forms $\Phi$, $\varphi$ with the same determinant, decide whether or not they are equivalent.}

\textit{Solution.} Find two reduced forms $F$, $f$ properly equivalent to the given forms $\Phi$, $\varphi$ respectively (Article 183). According as these are equivalent only properly, only improperly, in both ways, or in neither, the given forms will likewise be equivalent only properly, only improperly, in both ways, or in neither. Work out the period of either reduced form, \textit{e.g.}, that of $f$. If $F$ occurs in this period but the form associated with $F$ does not, we clearly have the \textit{first} case; if this associated form occurs but $F$ itself does not, the \textit{second}; if both occur, the \textit{third}; if neither, the \textit{fourth}.

\textit{Example.} Let the given forms be $(129,92,65)$, $(42,59,81)$, with determinant 79. We find the reduced forms $(10,7,-3)$, $(5,8,-3)$ properly equivalent to them. The period of the former is $(10,7,-3)$, $(-3,8,5)$, $(5,7,-6)$, $(-6,5,9)$, $(9,4.\,-7)$, $(-7,3,10)$. Since the form $(5,8,-3)$ itself does not occur in this period,
\SourceLeaf{190}
but its associated form $(-3,8,5)$ does, we conclude that the given forms are equivalent only improperly.

If all the reduced forms with a given determinant are divided into periods $P$, $Q$, $R$ etc., as above (Article 187, 5), and we choose any one form from each period (from $P$ $F$, from $Q$ $G$, from $R$ $H$ etc.), no two of these forms $F,G,H$ etc. can be properly equivalent. But every other form with the same determinant is properly equivalent to one of these, and to \textit{only one}. Clearly, then, \textit{all forms with this determinant can be divided into as many classes as there are periods}, by putting those properly equivalent to $F$ in the first class, those properly equivalent to $G$ in the second, and so on. In this way all forms in the same class will be properly equivalent, while forms in different classes cannot be properly equivalent. But we will not dwell on this subject here, since we will discuss it more fully below.

\Needspace{8\baselineskip}
\Article{196}
\textsc{Problem.} \textit{Given two properly equivalent forms $\Phi$, $\varphi$, find a proper transformation of one into the other.}

\textit{Solution.} By the method in Article 183 we can find two sequences of forms
\[
\Phi,\Phi',\Phi''\mathinner{\ldotp\ldotp\ldotp}\Phi^n\quad\text{and}\quad\varphi,\varphi',\varphi''\mathinner{\ldotp\ldotp\ldotp}\varphi^\nu
\]
such that each form is properly equivalent to its predecessor, and the last forms $\Phi^n$, $\varphi^\nu$ are reduced. Since $\Phi$, $\varphi$ are assumed properly equivalent, $\Phi^n$ must be contained in the period of $\varphi^\nu$. Let $\varphi^\nu=f$, and let its period as far as $\Phi^n$ be
\[
f,\ f',\ f''\mathinner{\ldotp\ldotp\ldotp} f^{m-1},\ \Phi^n
\]
so that $\Phi^n$ has index $m$ in this period. Denote the forms opposite to the associated forms of
\begin{center}
$\Phi,\ \Phi',\ \Phi''\mathinner{\ldotp\ldotp\ldotp}\Phi^n$\quad by $\Psi,\ \Psi',\ \Psi''\mathinner{\ldotp\ldotp\ldotp}\Psi^n$ respectively\HistoricalNote{*)}{Thus $\Psi$ is obtained from $\Phi$ by interchanging the first and last terms and giving the middle term the opposite sign; similarly for the others.}.
\end{center}
Then in the sequence
\SourceLeaf{191}
\[
\varphi,\ \varphi',\ \varphi''\mathinner{\ldotp\ldotp\ldotp} f,\ f',\ f''\mathinner{\ldotp\ldotp\ldotp} f^{m-1},\ \Psi^{n-1},\ \Psi^{n-2}\mathinner{\ldotp\ldotp\ldotp}\Psi,\ \Phi
\]
each form is a right neighbor of its predecessor. So by Article 177 we can find a proper transformation of the first form $\varphi$ into the last $\Phi$. This is easy to see for the other forms in the sequence; for $f^{m-1}$, $\Psi^{n-1}$, it is proved as follows. Let
\[
\begin{gathered}f^{m-1}=(g,h,i);\quad f^m\text{ or }\Phi^n=(g',h',i');\\ \Phi^{n-1}=(g'',h'',i'')\end{gathered}
\]
The form $(g',h',i')$ is a right neighbor both of $(g,h,i)$ and of $(g'',h'',i'')$; hence $i=g'=i''$ and $-h\equiv h'\equiv-h''\ (\mathrm{mod.}\ i\text{ or }g'\text{ or }i'')$. It follows that the form $(i'',-h'',g'')$, \textit{i.e.}, $\Psi^{n-1}$, is a right neighbor of $(g,h,i)$, \textit{i.e.}, $f^{m-1}$.

If the forms $\Phi$, $\varphi$ are improperly equivalent, $\varphi$ is properly equivalent to the form opposite to $\Phi$. So we can find a proper transformation of $\varphi$ into the form opposite to $\Phi$. If this is given by the substitution $\alpha,\beta,\gamma,\delta$, it is easy to see that $\varphi$ is improperly transformed into $\Phi$ itself by the substitution $\alpha,-\beta,\gamma,-\delta$.

This also shows that if the forms $\Phi$, $\varphi$ are both properly and improperly equivalent, we can find \textit{two} transformations, one proper and one improper.

\textit{Example.} We seek an improper transformation of the form $(129,92,65)$ into $(42,59,81)$, which we found in the preceding article to be improperly equivalent to it. We must therefore first find a proper transformation of $(129,92,65)$ into $(42,-59,81)$. For this purpose we work out the sequence of forms $(129,92,65)$, $(65,-27,10)$, $(10,7,-3)$, $(-3,8,5)$, $(5,22,81)$, $(81,59,42)$, $(42,-59,81)$. From this we derive the proper transformation $-47,56,73,-87$, which takes $(129,92,65)$ into $(42,-59,81)$; so the improper transformation $-47,\allowbreak-56,\allowbreak73,\allowbreak87$ takes it into $(42,59,81)$.
